% Class-compatibility preamble retained for a reader excerpt of
%   Sahar Diskin, Anna Geisler, "Perfect Matching in Product Graphs and in their
%   Random Subgraphs", arXiv:2404.14020v4, CC BY 4.0.
%
% The source pins its own publisher class (scrartcl); the delivery wrapper accepts
% only the classes article and amsart, so this file supplies the amsart baseline and
% then copies the author's own declaration block mechanically and unchanged
% (pinned-original-source.tex, source lines 30-57: the author's newtheorem
% declarations and the author's macro declarations).  The author's own package
% list of source lines 4-21 is reused only for the entries the retained text needs
% (enumitem, amsmath, amssymb, amsthm, mathtools, xcolor, geometry) and the
% cross-reference packages already needed by the author's own \Cref/\eqref usage
% (hyperref, cleveref); tikz, caption, subcaption, float, graphicx, wrapfig and
% abstract serve figures and the title block that this excerpt does not carry.
% No mathematical macro, statement, symbol or assumption is introduced, changed
% or removed by this file.
\documentclass[10pt,a4paper]{amsart}
\usepackage[shortlabels]{enumitem}
\usepackage{amsmath}
\usepackage{amssymb}
\usepackage{amsthm}
\usepackage{mathtools}
\usepackage{xcolor}
\usepackage[margin=1in]{geometry}
\usepackage{hyperref}
\usepackage{cleveref}
\hypersetup{
  colorlinks=true,
  linkcolor=blue,
  filecolor=magenta,
  urlcolor=cyan
}

\newtheorem{theorem}{Theorem}
\newtheorem{thm}{Theorem}[section]
\newtheorem{corollary}[thm]{Corollary}
\newtheorem{lemma}[thm]{Lemma}
\newtheorem{proposition}[thm]{Proposition}
\newtheorem{claim}[thm]{Claim}
\newtheorem{remark}[thm]{Remark}
\newtheorem{question}[thm]{Question}

\newcommand{\whp}{whp}
\newcommand{\defi}{\operatorname{def}}
\newcommand{\odd}{\operatorname{odd}}
\newcommand{\iso}{\operatorname{iso}}
\newcommand{\cN}{\mathcal{N}}
\newcommand{\cO}{\mathcal{O}}
\newcommand{\cE}{\mathcal{E}}
\newcommand{\cK}{\mathcal{K}}
\newcommand{\cP}{\mathcal{P}}
\newcommand{\cQ}{\mathcal{Q}}

\newcommand{\footremember}[2]{%
  \footnote{#2}%
  \newcounter{#1}%
  \setcounter{#1}{\value{footnote}}%
}

\newcommand{\ag}[1]{\textcolor{purple}{(Anna: #1)}}
\newcommand{\sd}[1]{\textcolor{red}{(Sahar: #1)}}
\crefname{originalequation}{original equation}{original equations}
\Crefname{originalequation}{Original equation}{Original equations}
\makeatletter
\newlabel{eq:expansion}{{6.2}{source}{}{sourcepointer-1}{}}
\newlabel{eq:expansion@cref}{{[originalequation][0][]6.2}{[1][0][]source}{}{}{}}
\makeatother
\crefname{originalequation}{original equation}{original equations}
\Crefname{originalequation}{Original equation}{Original equations}
\makeatletter
\newlabel{eq:f-star}{{6.1}{source}{}{sourcepointer-2}{}}
\newlabel{eq:f-star@cref}{{[originalequation][0][]6.1}{[1][0][]source}{}{}{}}
\makeatother
\title{A quantitative lower bound for the boundary function over the whole range of set sizes}
\author{Sahar Diskin, Anna Geisler}
\date{arxiv:2404.14020v4; CC BY 4.0}
\begin{document}
\maketitle
\noindent\textit{Editorial note.} All delivered mathematical text is the author's own, copied from the pinned author source sources/combinatorics-fl04-vspace-1-5cm-perfect-matching-in-p/pinned-original-source.tex; no statement or proof step was rewritten, reordered or removed.\par\medskip
\noindent\textit{Editorial note.} Preamble provenance: the source pins the publisher class scrartcl, which the delivery wrapper does not accept, so this card uses sources/combinatorics-fl04-vspace-1-5cm-perfect-matching-in-p/retained-preamble.tex, which supplies the amsart baseline and then carries the author's own newtheorem and macro declarations (pinned source lines 30-57) verbatim. No macro was invented and no mathematical symbol was changed.\par\medskip
\noindent\textit{Editorial note.} Bibliography provenance: the source carries no thebibliography environment and ends with bibliographystyle abbrv plus bibliography perc (author source lines 1691-1692), so the bibliography is resolved with BibTeX from the identical copy staged with the same arXiv version under sources/product-graph-matchings-2404.14020v4/perc.bib; no entry was edited and the author's own abbrv style is used. The counted proof itself contains no citation.\par\medskip
\noindent\textit{Editorial note.} Selection provenance: this card was opened under the file-based control contract (CONTROL.md), with the mechanical pre-card gate run on its statement and proof interval before the card was written and reporting the interval free.\par\medskip
\noindent\textit{Editorial note.} External references are navigation pointers only; every deliberately kept step of the author's argument stays inside the counted spans.\par\medskip
\noindent\textit{Allowed standard prerequisites (editorial).}\begin{itemize}
\item The boundary function f-star of the section and the two estimates on it coming from the expansion inequalities. Set up in the author's section on the random graph process, with the author's displays labelled eq:expansion and eq:f-star (source lines 795-810). Supplies the two lower bounds for the boundary function that the proof evaluates in each range of the size parameter.
\end{itemize}\medskip
\ifdefined\hypertarget\hypertarget{sourcepointer-1}{}\fi
\noindent\textit{Source reference.} Original equation 6.2: sources/combinatorics-fl04-vspace-1-5cm-perfect-matching-in-p/pinned-original-source.tex, source lines 807--807. Navigation pointer only: the counted proof uses the author's displayed expansion inequality as one of its two inputs; the display is not reproduced in this card.\par\medskip
\ifdefined\hypertarget\hypertarget{sourcepointer-2}{}\fi
\noindent\textit{Source reference.} Original equation 6.1: sources/combinatorics-fl04-vspace-1-5cm-perfect-matching-in-p/pinned-original-source.tex, source lines 801--801. Navigation pointer only: the counted proof uses the author's displayed two-term estimate on the boundary function; the display is not reproduced in this card.\par\medskip
\begin{lemma}\label{lem:strong-boundary}
For all sufficiently large $d$ and $3\le k\le2u_1$,
\begin{align}
 f^*(k)&\ge2d+15 (M-1) k \ln d.\label{eq:strong-3d}
\end{align}
\end{lemma}

\begin{proof}
For $k=3$, \eqref{eq:expansion} together with the first term of \eqref{eq:f-star} give
\[
 f^*(k)\ge 3d- 3(M-1) \log_M(3) \geq 2d + 45(M-1) \ln d,
\]
as claimed. For $4 \leq k \le d^3$ the same implies
\[
f^*(k)\ge kd - 3k(M-1) \log_M(d) \geq 3d + 15k (M-1) \ln d.
\]
If $d^3<k\le2u_1$, then \eqref{eq:expansion} together with the second term of \eqref{eq:f-star} give
\[
 f^*(k)\ge\frac{ek}{M}\left(\frac{M^3}{p_-}\ln d-\ln 2\right).
\]
Using that $p_-\le1-1/M$
%since n \leq M^d.
, we have $\frac{eM^2}{p_-} \geq 16(M-1)+1$. This yields
\[
 f^*(k)\ge 16(M-1) k \ln d + k \ln d - \frac{ek}{M} \ln 2 \geq 3d + 15(M-1)k \ln d,
\]
where the last inequality uses that $k>d^3$.
\end{proof}

\bibliographystyle{plain}
\bibliography{perc}
\end{document}
